% qa_notes_ko.tex (2026-10-22): 연구 질문 · 답변 · 방향 정리 문서. 빌드: bash paper/manuscript/qa_notes/build.sh -> qa_notes_ko.pdf (deliverables/paper/ 로 복사)
\documentclass[10pt,a4paper]{article}
\usepackage{xeCJK}\setCJKmainfont{Noto Serif CJK KR}\xeCJKsetup{CJKspace=true}
\usepackage[margin=22mm]{geometry}\usepackage{booktabs,longtable,array,tabularx}\usepackage[hidelinks]{hyperref}
\setlength{\parskip}{3pt}\setlength{\parindent}{0pt}\renewcommand{\arraystretch}{1.12}
\title{\large 물리 잔차 ALT 연구: 질문 · 답변 · 방향 정리}
\author{\normalsize Polar\_Bigdata 세션 기록(원고 v15.2 기준)}\date{\normalsize 2026-10-22, 최신 갱신}
\begin{document}\maketitle

\section{문서의 목적과 관리 규칙}
이 문서는 원고 검토 중에 나온 질문(2026-10-21 · 10-22)과 그 답, 그리고 연구의 목적 · 결과 · 결론에 대한 정성적 서술을 한 곳에 모은다. 원고 본문이 아니라 저자 내부 참고용이며, 질문이 더해질 때마다 3절에 항목을 추가하고 8절에 변경 이력을 남긴다. 수치는 모두 저장된 결과 표(\texttt{data/processed/v5/runs/core\_v6}, \texttt{ladder\_v6})와 그림 source\_data 에서 가져왔고, 본문 매크로와 같은 값이다. 소스: \texttt{paper/manuscript/qa\_notes/qa\_notes\_ko.tex}.

\section{연구의 목적 · 명분 · 의의}
\textbf{문제.} ALT 관측은 드물고 군집되어 있어 지역 지도는 관측을 확장하는 모델에 의존한다. Stefan 모델은 전이되지만 지역당 계수 하나로 토양 · 지형 차이를 담지 못하고, 직접 ML 은 학습 지점을 맞추지만 학습 지점에서 떨어진 곳과 새 지역에서 정확도가 떨어지며, 공개 격자 제품은 자기 지점 집합 안에서만 검증되어 있다.

\textbf{접근.} 재보정 Stefan 모델을 기준으로 두고 그 잔차를 기후 · 지형 · 토양 23개 입력으로부터 gradient boosting 이 학습하는 물리 잔차 워크플로. 모든 라벨 · 입력 · 예측 · 점수가 '지점 하나 · 명시된 한 해'를 가리키는 지점-연도 설계, 다섯 검증 설계(무작위 → 지역 제외), 공개 제품 5종을 같은 지점-연도에서 채점.

\textbf{주장(원고 제목 · 초록).} (i) 방법의 순위는 검증 설계에 따라 바뀌며 직접 ML 은 지역이나 라벨 대부분을 빼면 가장 많이 잃는다. (ii) 잔차 학습기는 토양 · 지형 입력이 있을 때 물리를 개선하고, 새 지역에서는 라벨 25개부터 이득을 준다. (iii) 공개 제품의 지역 차이는 학습 지점 분포로 설명되며, 같은 지점-연도에서 물리 잔차 모델이 제품보다 관측에 가깝다. (iv) 적용 가능 영역과 예측 구간이 딸린 2011--2024년 연도별 1 km 지도 세 지역.

\textbf{의의.} 기법 자체(잔차 학습)는 새롭지 않다. 기여는 ALT 에 대해 물리 기준 잔차 학습을 지점-연도 단위로, 다섯 검증 설계와 라벨 수 · 전이 실험으로, 공개 제품과 같은 조건에서 처음 체계적으로 시험한 평가 설계와, 관측이 드문 지역에 적용할 수 있는 절차(물리 출발점 + 소수 지역 라벨)에 있다.

\section{질문과 답변}

\subsection*{Q1 (10-21). 알래스카 2013 공통 셀에서 우리 지도의 상관(0.34)이 제품(0.53 · 0.57)보다 낮은데 중요한가, 2013 을 쓸 이유가 있나}
2013 은 알래스카에서 관측 셀이 가장 많은 해이므로 선택 기준(지역별 관측 셀 최다 해)에 따른 것이다. 상관이 낮은 이유는 알래스카 2013 공통 셀 126개의 관측 범위가 좁고(대부분 40--80 cm) 셀 안 산포(합동 표준편차 13 cm)가 크기 때문이며, RMSE 로 보면 우리 14.9 cm 대 Wei 13.4 · Liu 12.9 cm 로 차이가 작고 25 km 밖에서는 9.1 대 10.9 · 8.5 cm 이다. 이득 부재는 2013 이나 가중치 $\lambda$ 의 문제가 아니다(평균 입력 +0.33, 연도 입력 $-0.15$, $\lambda$ 0.25/1 에서 $-0.15$/+0.30). 알래스카에서는 라벨이 15,540개로 조밀해 재보정 Stefan 이 이미 지역 수준을 맞추고, 남는 오차의 큰 부분이 1 km 안의 산포여서 어느 방법도 줄이지 못한다(보충 표 S6). 원고는 이를 '알래스카 이득 부재' 로 그대로 적는다.

\subsection*{Q2 (10-21). 1 km 셀 단위 라벨 재집계는 왜 필요하고 전면 재실행이 필요한가}
한 1 km 셀에 여러 탐침 지점이 들어가므로(알래스카 중앙값 여러 지점) 지점-연도 점수는 셀 안 산포를 모델 오차로 센다. 같은 겹 밖 예측을 셀 · 연도별로 평균해 다시 채점하면(재실행 없이 재집계만) 지도 해상도에서의 오차를 얻는다. 결과(보충 표 S6): 물리 잔차 대 Stefan 의 이득은 두 단위에서 유지되고, 알래스카에서 직접 ML 의 지점-연도 우위(17.3 대 18.4)는 셀-연도에서 뒤집힌다(24.7 대 24.0); 지역 제외에서는 두 단위 모두 직접 ML 이 가장 부정확하다(알래스카 23.9 대 23.4, 레나 37.4 대 22.2). 전면 재실행은 필요하지 않다.

\subsection*{Q3 (10-21 · 10-22). 지도를 2011--2024년 평균으로 그리는 선행 연구 근거는 무엇이고, PDF 에 어떻게 기록되어 있나}
확인한 PDF(\texttt{references/}) 와 캡션 원문의 요지는 다음과 같다(상세 조사: \texttt{docs/research/2026-10-21/map\_and\_figure\_practice.md}).
\begin{small}\begin{tabularx}{\textwidth}{@{}>{\raggedright\arraybackslash}p{0.2\textwidth}>{\raggedright\arraybackslash}p{0.34\textwidth}X@{}}
\toprule 논문 · PDF & 본문 지도 & 캡션 원문(요지) \\ \midrule
Liu et al. 2024, ERL 19, 014020 (\texttt{08\_recent\_alt\_2022\_2026/liu2024\_widespread\_alt\_deepening\_2003\_2020\_erl.pdf}) & 2003--2020 평균 지도 + 추세 지도; 연도별 지도 없음 & Figure 4 ``Spatial patterns of (a) mean, (b) standard deviation, (c) coefficient of variation \ldots''(2003--2020); Figure 5 ``\ldots (a) mean \ldots'' of the trend; Figure 3 ``site-level comparison \ldots of mean ALT'' \\
Wei et al. 2026, ESSD Discussions (\texttt{08\_recent\_alt\_2022\_2026/wei2026\_nh\_alt\_1km\_2000\_2024\_essdd.pdf}) & 선택 연도 소형 다중(a--f) + 면적 비율; 제품 비교는 2000--2020 평균 & Figure 4 ``Spatiotemporal distribution of ALT (a--f) and the proportion of ALT area at different levels (g) from 2000 to 2024''; Figure 7 ``Spatial comparison of the reconstructed ALT dataset with existing 1 km hemispheric products'' \\
Ran et al. 2022, ESSD 14, 865 (\texttt{01\_benchmark/ran2022\_panarctic.pdf}) & 2000--2016 평균 & Figure 4 ``Distribution of the average active layer thickness (ALT) in the Northern Hemisphere for the period of 2000--2016'' \\
Aalto et al. 2018, GRL 45, 4889 (\texttt{12\_manuscript\_refs/aalto2018\_circumarctic\_ground\_temp\_alt.pdf}) & 기간 평균(2000--2014, 2041--2060)과 차이 & Figure 4 ``Forecasted mean annual ground temperatures (MAGT) and active layer thickness (ALT) across study domain for two periods and their differences'' \\
Yi et al. 2018, The Cryosphere 12, 145 (조사 문서 11번, pdftotext 캡션) & 2001--2015 평균 ALT + 추세 지도 & Fig 3a 다년 평균 1 km ALT, Fig 5a 2001--2015 추세 \\
Obu et al. 2019, Earth-Sci. Rev. 193, 299 (\texttt{02\_alt\_dl\_mapping/obu2019\_ttop\_nh.pdf}) & 2000--2016 평균 MAGT & Fig 3 다년 평균 MAGT(ALT 지도 없음) \\
ESA CCI Permafrost v4 PUG (\texttt{01\_benchmark/esa\_cci\_permafrost\_v4.pdf}) & 연도별 제품; 설명서 예시는 단일 연도 & ALT 2015 예시 지도 \\
\bottomrule\end{tabularx}\end{small}
결론: 본문 주 지도를 다년 평균으로 두는 것이 ALT · 영구동토 ML 지도 논문의 지배적 관행이고(위 6편 중 5편), 연도별 지도는 선택 연도 소형 다중(Wei)이나 단일 연도 예시(CCI)로만 쓰인다. 우리 그림 7 은 평균 지도(a--c)에 연도별 편차 14장(f) · TDD 응답(d) · 추세(e)를 더하고, 그림 6 은 특정 연도의 관측 셀과 제품을 같은 해에서 비교하므로 관행 안에 있으면서 극지연의 '연도 불명' 지적을 해소한다. 권고: 그림 7a--c · 보충 그림 S8 은 '14장 연도별 지도의 평균' 으로 두고 캡션에 명시하며(현재 반영), 특정 연도 비교는 그림 6 · 보충 그림 S5 가 맡는다. 제품 비교까지 단일 연도로 한 예는 선행 연구에 없으므로, 이 점은 원고에서 '처음' 으로 적되 과장하지 않는다.

\subsection*{Q4 (10-22). 보충 방법 S3 의 제품들은 다른 지역 예측을 염두에 둔 것인가, 공개되어 사용이 허용되는가}
\begin{itemize}\setlength{\itemsep}{1pt}
\item Wei et al. 2026 (ESSD Discussions, 저자 CC BY 4.0; 자료 Zenodo 21667583 v2, CC BY 4.0): ``A 1 km resolution dataset of Northern Hemisphere permafrost active layer thickness \ldots from 2000 to 2024''; 연구 영역은 ``the extensive permafrost domain of the Northern Hemisphere \ldots including Siberia, Alaska, and the Canadian Arctic''. 반구 전체 제품이며 임의 지역 사용이 전제이다. 토론 단계 원고이므로 인용 시 ESSDD 로 표기한다.
\item Liu et al. 2024 (ERL, open access CC BY 4.0; 자료 Zenodo 10070610, CC BY 4.0): ``northern permafrost regions (NPRs) from 2003 to 2020''. 반구 제품.
\item ESA CCI Permafrost ALT v4 · v5 (CEDA, ESA CCI 자료 정책: 출처 표기 조건의 자유 이용): 범북극 연도별 1 km 제품(1997--2021).
\item Aalto et al. 2018 (GRL; 자료 Zenodo 5003804, CC0): ``circum-Arctic'' 영역(북위 30도 이상 육지), 2000--2014 기준기 값.
\item Yi \& Kimball 2019 (ORNL DAAC 1760, NASA EOSDIS 공개): 알래스카 한정이므로 알래스카에만 썼다.
\end{itemize}
따라서 다섯 제품 모두 공개 약관으로 재사용이 허용되고, Yi \& Kimball 을 뺀 넷은 반구 전체를 덮는 제품이라 우리 세 지역에서의 채점은 제품의 의도된 사용 범위 안이다. 다만 어느 제품도 '학습 지점을 주지 않은 지역의 특정 해 관측' 으로 검증된 적은 없고, 이것이 그림 5 · 6 의 명분이다. 원고 보충 방법 S3 과 보충 표 S1 에 약관 · 저장소를 적었다.

\subsection*{Q5 (10-22). 라벨 수를 늘리면서 어떤 자료를 어떻게 늘렸나}
\textbf{지역 안 사다리(보충 그림 S4).} 지역의 지점-연도를 0.5도 공간 블록 5겹으로 나눈 뒤, 겹마다 학습 겹(다른 4겹)의 지점-연도에서 $n$개(25, 50, 100, 250, 500, 1000, 2500, 전량)를 무작위로(균등, 비복원, 씨앗 0 · 1 · 2) 뽑아 모델을 학습하고, 제외한 겹의 모든 지점-연도로 채점한다. 기준은 같은 $n$개로 재보정한 Stefan 모델이고, 차이는 5겹을 합친 뒤 블록 부트스트랩으로 구간을 구한다. 추출은 블록 · 연도로 층화하지 않으므로 한 추출이 공간적으로 몰릴 수 있고, 3회 평균이 이를 완화한다. \textbf{전이 사다리(그림 4c).} 원천 = 다른 두 지역의 전체 지점-연도, 대상 라벨 = 대상 지역 학습 겹에서 무작위 $n$개(0, 25, 50, 100, 250, 500), 시험 = 대상 지역의 제외 겹. Stefan 은 $n \ge 3$ 이면 대상 라벨로 재보정, 아니면 원천 계수. 코드: \texttt{scripts/3\_deep\_learning/v6\_label\_ladder.py}.

\subsection*{Q6 (10-22). 입력 연도가 섞이거나 평균 입력이 쓰인 실험은 없는가(변인 통제 감사)}
코드 경로를 따라 확인한 결과이다.
\begin{itemize}\setlength{\itemsep}{1pt}
\item 라벨: 지점-연도(관측 위치 · 관측 연도). 한 지점의 한 해 여러 측정은 평균, 평균 관측일 보존.
\item 기후 입력: \texttt{testbed\_v5.with\_forcing(\ldots, "year\_obs")} 가 관측 연도의 월 자료에서 관측일까지 누적한 TDD 를 쓴다(날짜 없는 CALM 은 9월 1일). 나머지 7개 기후 변수도 관측 연도 값. 전년 지연 변수(민감도)는 관측 전년.
\item 정적 입력: 지형 · 토양(시간 불변). MODIS(민감도 묶음)는 \texttt{build\_v5\_modis\_cell\_year.py} 가 셀 · 연도별로 만든 값을 쓴다.
\item 지도: \texttt{v5\_map\_predict.py} 가 해마다 그 해의 기후 입력(TDD 는 9월 1일까지)으로 예측하고, 2011--2024년 평균 지도는 14장 연도별 지도의 평균이다(평균 입력으로 만든 지도가 아님). 적용 가능 영역의 비유사도 지수만 2011--2024년 평균 입력으로 계산한다(방법에 기술).
\item 제품: 관측 연도의 제품 값(Aalto 는 단일 기준기 값).
\item 평균 기후 입력은 그림 3a 의 민감도 비교에만 쓰이고, 보충 표 S2 가 분석별 연도 · 입력을 나열한다.
\end{itemize}
따라서 '연도 혼재' 실험은 없다. 설계상 주의점 하나: 지도는 9월 1일 TDD(계절 말) 기준이고 라벨 채점은 관측일 기준이므로, 지도와 관측을 직접 비교하는 그림 6 · 보충 그림 S10 은 이 차이를 캡션에 밝힌다.

\subsection*{Q7 (10-22). 직접 ML 만으로 물리 제품이나 Stefan 보다 좋게 나올 수 있나, 어느 해의 예측인가}
가능하다. 지역 안 공간 블록 검증에서 직접 ML 은 그 지역의 라벨 수천 개와 지형 · 토양 입력으로 학습하므로 계수 하나인 Stefan 보다 지점 차이를 잘 맞추고(알래스카 17.3 대 18.6 cm), 공개 제품은 우리 라벨로 학습하지 않았거나(Wei 는 CALM · GTN-P 201개 지점) 관측으로 학습하지 않은(CCI) 모델이므로 우리 지역에서 더 나쁠 수 있다. 모든 수치는 2010--2024년 지점-연도 전부를 각자 자기 해의 입력과 제품 값으로 채점한 것이다(연도별 결과는 보충 그림 S5). 그러나 이 우위는 시험 지점이 학습 지점 사이에 끼어 있는 설계에서만 나타나고, 지역을 빼면(레나 39.3 cm) 또는 라벨이 적으면(알래스카 · 캐나다 25--250개에서 Stefan 보다 나쁨) 사라진다. 원고는 이 사실을 숨기지 않되 순서를 바꿔, 물리 잔차 모델의 이득과 안정성을 먼저 적고 직접 ML 의 전량 라벨 우위는 조건부 사실로 뒤에 둔다(보충 표 S4 · S6 캡션의 '읽는 법').

\subsection*{Q8 (10-22). 지점-연도와 셀-연도의 차이}
지점-연도 = 관측 위치 하나(좌표 0.0001도 반올림)의 한 해. 셀-연도 = 한 1 km 지도 셀 안 모든 지점의 관측과 예측을 한 해 단위로 평균한 것. 알래스카 15,540 지점-연도가 910 셀-연도가 되는 것은 ABoVE 탐침 측선의 여러 지점이 한 셀에 들어가기 때문이다. 보충 방법 S1 과 보충 표 S6 캡션에 정의를 넣었다.

\subsection*{Q9 (10-22). 그림에서 'Wei' 라고 적어도 되나}
그림 안 짧은 표기로 'Wei 2026' 은 흔하지만 'Wei et al. 2026' 이 안전하다. 모든 그림 · 표 · 보충의 제품 표기를 'Wei et al. 2026', 'Liu et al. 2024', 'Aalto et al. 2018' 로 통일하였다(색표 공용 모듈과 각 스크립트).

\subsection*{Q10 (10-22). 옛 보충 그림 S5a 에서 캐나다 직접 ML 의 구간이 길었던 이유}
연도별 입력 대 평균 입력의 RMSE 차이에 대한 95\% 블록 부트스트랩 구간인데, 캐나다는 블록 26개뿐이고 직접 ML 의 블록별 차이가 블록마다 부호가 달라 구간이 넓다(그림 3a 에서도 캐나다 직접 ML 이 가장 넓다). 판정은 '차이 미해결' 이며 모델이 나쁘다는 뜻이 아니다. 옛 S5a 는 본문 그림 3a 와 중복이라 보충에서 뺐다.

\subsection*{Q11 (10-22). 보충 그림 S10 에 레나 제품 값이 없고 캐나다 패널이 없던 이유}
레나 델타의 최다 관측 지점(사모일로프 부근)은 Wei · Liu 래스터에 값이 없는 셀(수역 · 삼각주 마스크)이라 제품 선을 그릴 수 없다. 캐나다는 옛 스크립트가 두 지역만 골랐기 때문이며, v15.2 에서 세 지역으로 바꾸고 관측을 선으로 잇고, 지역별 가장 따뜻한 해 · 추운 해를 음영으로 표시해 '튀는 해' 가 더운 계절(알래스카 2019, 레나 2021)임을 읽게 하였다. 캐나다 지점은 관측이 모델보다 약 20 cm 깊은 수준 차이를 보이는데, 이는 CALM 격자 평균(원고 보충 표 S7 의 캐나다 CALM 편향 $-23$ cm)과 같은 현상이다.

\subsection*{Q12 (10-21). 왜 Stefan 인가, 왜 CatBoost 인가}
Stefan: 매개변수 하나라 소수 관측으로 재보정 · 전이되고, 입력이 연도별 TDD 라 연도 응답을 가지며, Kudryavtsev 는 블록 CV 에서 더 나쁘다(21.7/20.4/25.4 대 Stefan 18.6/21.0/30.6 cm; 레나 · 캐나다에서 우세하나 지역 제외에서 25.1/21.8/35.4 로 불안정). CatBoost: 학습기 5종(CatBoost · RF · MLP · TabM · FT-T) 비교에서 물리 잔차 구조의 이득이 가장 안정적이고(보충 그림 S2b), 표 자료에서 gradient boosting 이 신경망과 대등하거나 낫다는 선행 결과(Grinsztajn 2022)와 맞다.

\section{결과 정리와 정성적 결론}
\begin{small}\begin{tabular}{@{}lrrrrrr@{}}\toprule
RMSE (cm), 지점-연도 & \multicolumn{2}{c}{알래스카} & \multicolumn{2}{c}{레나 델타} & \multicolumn{2}{c}{캐나다(북서부)} \\
 & 블록 CV & 지역 제외 & 블록 CV & 지역 제외 & 블록 CV & 지역 제외 \\ \midrule
Stefan(재보정 / 원천) & 18.6 & 19.4 & 21.0 & 21.5 & 30.6 & 27.8 \\
물리 + 잔차 ML & 18.4 & 17.5 & 17.4 & 23.5 & 26.9 & 24.2 \\
직접 ML & 17.3 & 20.6 & 15.9 & 39.3 & 26.1 & 23.9 \\ \midrule
전이, 대상 라벨 0 / 25 / 500 (물리 잔차) & \multicolumn{2}{c}{17.8 / 17.3 / 17.5} & \multicolumn{2}{c}{23.8 / 19.9 / 18.5} & \multicolumn{2}{c}{24.3 / 26.3 / 26.1} \\
전이, 대상 라벨 0 / 25 / 500 (직접 ML) & \multicolumn{2}{c}{21.5 / 19.5 / 18.5} & \multicolumn{2}{c}{39.5 / 21.3 / 17.5} & \multicolumn{2}{c}{24.2 / 23.8 / 22.8} \\
지역 안 라벨 25개, Stefan 대비 $\Delta$ (물리 잔차 / 직접) & \multicolumn{2}{c}{$-0.9$ / $+2.8$} & \multicolumn{2}{c}{$-0.8$ / $-0.8$} & \multicolumn{2}{c}{$-0.3$ / $+1.2$} \\
공통 지점-연도: 물리 잔차 / Wei / Liu & \multicolumn{2}{c}{17.1 / 17.9 / 18.8} & \multicolumn{2}{c}{21.1 / 34.6 / 33.0} & \multicolumn{2}{c}{25.8 / 27.0 / 27.9} \\
관측에 더 가까운 셀 비율(대 Wei / 대 Liu) & \multicolumn{2}{c}{52 / 50 \%} & \multicolumn{2}{c}{82 / 85 \%} & \multicolumn{2}{c}{53 / 42 \%} \\
\bottomrule\end{tabular}\end{small}

\textbf{Stefan 대비.} 레나 · 캐나다에서 3.6 · 3.7 cm(약 17 · 12 \%) 개선, 알래스카는 개선 없음. 이득은 토양 · 지형 입력에서 나오고(기후만으로는 없음), 잔차는 Stefan 항에 없는 정보를 학습한다(SHAP).

\textbf{직접 ML 대비.} 전량 라벨 · 지역 안 설계에서는 직접 ML 이 0.8--1.6 cm 더 정확하다(셀-연도에서는 알래스카가 뒤집힘). 그러나 지역을 빼면 직접 ML 은 알래스카(20.6 대 17.5)와 레나(39.3 대 23.5)에서 크게 나빠지고, 라벨 25--250개에서는 알래스카 · 캐나다에서 Stefan 보다 나쁘며, 전이에서 대상 라벨이 적을 때 알래스카 · 레나에서 물리 잔차 모델이 앞선다. 캐나다 전이에서는 직접 ML 이 모든 $n$ 에서 1.5--3 cm 더 정확한데, 이는 알래스카의 ABoVE 탐침 자료가 캐나다와 같은 프로그램 · 규약이라 원천 라벨이 대상을 잘 대표하기 때문으로 해석한다. 즉 물리 잔차 모델의 장점은 '가장 정확' 이 아니라 '설계 · 지역 · 라벨 수에 걸쳐 가장 안정적이고 Stefan 보다 나빠지지 않음' 이다.

\textbf{제품 대비.} 같은 지점-연도에서 세 지역 모두 물리 잔차 모델이 Wei · Liu 보다 RMSE 가 낮고, 레나에서 차이가 크다(제품은 관측과 무관한 좁은 범위를 예측). 알래스카 · 캐나다에서는 차이가 1--2 cm 로 작아 '약간 나은 수준' 이 정직한 표현이다. 지역을 빼고도(지역 제외 RMSE 17.5 / 23.5 / 24.2) 제품(17.9 / 34.6 / 27.0)과 같거나 낫다는 점이 전이성 주장의 근거이다.

\section{선행 연구 대비 차별점}
(1) 지점-연도 설계: 라벨 · 입력 · 예측 · 점수가 한 해를 가리키며 입력은 관측일까지 누적(제품은 연도 평균 변수 또는 다년 평균 라벨). (2) 다섯 검증 설계를 같은 자료에 적용해 순위 변화를 보임(제품 논문은 무작위 또는 지점 제외 분할 하나). (3) 라벨 수 · 전이 사다리로 '새 지역에 몇 개가 필요한가' 를 수치로 제시. (4) 공개 제품 5종을 같은 지점 · 같은 해에서 채점하고 학습 지점 거리(25 km)로 나누어 비교. (5) 적용 가능 영역과 예측 구간을 지도에 함께 제공. 기법(잔차 학습, gradient boosting)은 Willard 2022 · Read 2019 의 틀 안에 있고, Gautam 2025(알래스카, Stefan 대 RF)를 세 지역 · 다섯 설계 · 제품 비교로 확장한 위치이다.

\section{기대효과 · 적용 가능성 · 한계}
\textbf{기대효과.} 관측이 드문 지역에서 (a) 원천 계수 Stefan 으로 시작하고 (b) 지역 관측 25개 안팎으로 재보정 · 잔차 보정하며 (c) 적용 가능 영역 밖은 쓰지 않는 절차. 연도별 지도는 기후의 해별 변동을 담으므로 지반 안정성 · 수문 · 탄소 연구에서 특정 해의 융해 깊이를 요구하는 응용(예: 더운 해의 융해 심화)에 쓸 수 있다. \textbf{결합 가능성.} 보충 그림 S7 은 잔차 학습기를 제품(Wei) 기준에도 붙일 수 있음을 보이며(알래스카 19.6 → 18.7 cm), Stefan 기준이 여전히 더 정확하다. 물리 모델 두 개의 앙상블(Stefan + Kudryavtsev)과 제품 결합은 후속 과제로 적는다. \textbf{한계.} 알래스카에서 이득이 없고, 1 km 안 산포(13 cm)는 어느 방법도 줄이지 못하며, 캐나다 전이는 직접 ML 이 더 나아 '물리 기준이 항상 낫다' 고 말할 수 없다. 추세는 모델 응답이지 관측 추세가 아니다.

\section{냉정한 평가(2026-10-22)}
초록이 제시한 세 비교 가운데 Stefan 대비 개선(두 지역)과 제품 대비 우위(레나 명확, 나머지 약간)는 달성했고, 직접 ML 대비는 '조건부' 다. 목표를 '모든 설정에서 가장 정확한 모델' 로 잡으면 달성하지 못했고, '물리의 전이성을 지키면서 ML 의 정확도에 가까워지는 모델' 로 잡으면 달성했다. 원고는 후자로 명분을 세운다: 제목 · 초록 · 서론 · 결론을 '안정성 · 전이성 · 새 지역 절차' 로 맞추고, 직접 ML 의 전량 라벨 우위는 사실로 적되 그 조건(지점 끼워 넣기 설계)과 붕괴(지역 제외 · 소수 라벨)를 같은 문장에 둔다. 사용자가 든 '특정 라벨 조건에서 가장 뛰어난 지점이 있고, ML 과 물리 잔차 ML 의 혼용 · 제품 결합의 가능성을 연 것' 이라는 판단에 동의하며, 혼용 규칙(예: 라벨 수 · 적용 가능 영역에 따라 직접 ML 과 물리 잔차를 고르는 결정 규칙)은 수치가 더 필요해 고찰의 '향후' 로 둔다.

\section{변경 이력}
\begin{itemize}\setlength{\itemsep}{1pt}
\item 2026-10-22: 문서 신설. Q1--Q12, 결과 표, 평가 수록. 원고 v15.2(보충 재구성 · 초록 · 서론 수정)와 동기.
\end{itemize}
\end{document}
